\clearpage
\VerdierPhysicalPageAnchor{0158}%
\begingroup\small\setlength{\parskip}{0pt}\setlength{\emergencystretch}{1.5em}
\VerdierSemanticAnchor{ch3-3-2-5-kernel-triangle}
$(Z,X,Y',u,f'g,v)$ を $\mathcal D$ の完全三角とする。図式
\[
\begin{tikzcd}[column sep=5em,row sep=2.1em]
h(Z)\arrow[r,"h(fu)"]\arrow[d,"h(u)"']&h(Y)\arrow[d,"\operatorname{id}_{h(Y)}"]\\
h(X)\arrow[r,"h(f)"]\arrow[d,"h(g)"']&h(Y)\arrow[d,"h(g')"]\\
h(X')\arrow[r,"h(f')"']&h(Y')
\end{tikzcd}
\]
の横射の像をとると
\VerdierSemanticAnchor{ch3-3-2-5-1}
\[
\operatorname{Im}(fu)\xrightarrow{\beta}V\xrightarrow{\alpha}V'.
\tag{3.2.5.1}
\]
を得る。
\VerdierSemanticAnchor{ch3-3-2-5-kernel-proof}
射 $\operatorname{Im}(fu)\to V$ が $\alpha$ の核であることを示す。
まず $\beta$ はモノ射である。実際、$\operatorname{Im}(fu)$ と $V$ は
$h(Y)$ の部分対象である。次に $\alpha\beta=0$ である。標準射
$h(Z)\to\operatorname{Im}(fu)$ はエピ射、標準射
$V'\to h(Y')$ はモノ射であり、$h(f'gu)$ は完全三角の連続する
二射の合成の $h$ による像だから零である (1.2.2)。
残るは $V$ での完全性である。$W$ を $\mathcal D$ の対象とする。
任意の射 $\sigma\colon h(W)\to V$ は、$h(W)$ が $\mathcal D^\wedge$ で射影的なので、
$h(W)$ から $h(X)$ への射、すなわち射 $m\colon W\to X$ から得られる。
$\alpha$ との合成が零であることと $f'gm=0$ は同値である。これは
$V'$ が $h(Y')$ の部分対象だからである。(1.2.1) により
$un=m$ なる $n\colon W\to Z$ が存在する。よって $\sigma$ は
$\operatorname{Im}(fu)$ を経由する。

\VerdierSemanticAnchor{ch3-3-2-5-cokernel-setup}
次に $\mathcal D^\wedge$ における $\alpha$ の余核が
$\mathcal D^\wedge$ の許容対象であることを示す。
$(X,Y',Z',g'f,u',v')$ を $\mathcal D$ の完全三角とする。図式
\[
\begin{tikzcd}[column sep=5em,row sep=2em]
h(X)\arrow[r,"h(f)"]\arrow[d,"h(g)"']&h(Y)\arrow[d,"h(g')"]\\
h(X')\arrow[r,"h(f')"]\arrow[d,"\operatorname{id}_{h(X')}"']&h(Y')\arrow[d,"h(u')"]\\
h(X')\arrow[r,"h(u'f')"']&h(Z')
\end{tikzcd}
\]
\endgroup
